 \noindent
 Let $\kk$ be any field of characteristic zero. Write $\KK=\kk(X,Y)$ for the skew-field of rational functions in non-commuting variables $X$ and $Y$. Intuitively, writing $\pi:\kk(X,Y)\to\kk(x,y)$ for the commutative specialization, we may formally invert any element $W\in\KK$ for which $\pi(W)\ne0$; this idea has been made precise in~\cite{usnich1} by considering iterated localizations of the free algebra $\kk\langle X,Y\rangle$.

 For any nonzero polynomial $P\in\kk[z]$, consider the following $\kk$-linear automorphism of $\KK$:
 \[F_P:\begin{cases} X\mapsto XYX^{-1}\\ Y\mapsto P(Y)X^{-1}.\end{cases}\]
 We remark for later use that the element $Q\coloneqq XYX^{-1}Y^{-1}$ is fixed by $F_P$ for any nonzero polynomial $P$.
 Also note that $F_P^{-1}$ is given by $X\mapsto P(X)Y^{-1}$ and $Y\mapsto YXY^{-1}$.

 Fix nonzero monic polynomials $P_1,P_2\in\kk[z]$ such that $P_1(0)=1=P_2(0)$, say
 \[P_i(z)=p_{i,0}+p_{i,1}z+\cdots+p_{i,d_i-1}z^{d_i-1}+p_{i,d_i}z^{d_i}\]
 with $p_{i,0}=p_{i,d_i}=1$ for $i=1,2$. Set $\AA_+=\ZZ_{\ge0}[p_{1,i},p_{2,j}:0\le i\le d_1,0\le j\le d_2]$ and call this the \emph{pseudo-positive semiring} associated to $P_1$ and $P_2$. 

 We will write $\bar{P}_1(z)\coloneqq z^{d_1}P_1(z^{-1})$ and $\bar{P}_2(z)\coloneqq z^{d_2}P_2(z^{-1})$ for the polynomials obtained from $P_1$ and $P_2$ by reversing the order of the coefficients. Note that these are again polynomials of the same form. For notational convenience, for $k\in\ZZ$ we define 
 \begin{equation}
 \label{eq:reversed polynomials}
 P_k=p_{k,0}+p_{k,1}z+\cdots+p_{k,d_k-1}z^{d_k-1}+p_{k,d_k}z^{d_k}=\begin{cases}\bar{P}_2 & \text{if $k\equiv 0\mod 4$;}\\P_1 & \text{if $k\equiv 1\mod 4$;}\\P_2 & \text{if $k\equiv 2\mod 4$;}\\\bar{P}_1 & \text{if $k\equiv 3\mod 4$.}\end{cases}
 \end{equation}
 Here we use the notation $d_k\coloneqq d_1$ if $k$ is odd and $d_k\coloneqq d_2$ if $k$ is even.

 Set $X_0\coloneqq X$ and $Y_0\coloneqq Y$.
 \begin{maintheorem}
 \label{th:main}
 For any $m\ge0$, the elements $ X_m, Y_m\in\KK$ given by
 \begin{equation}
 \label{eq:non-commutative initial cluster mutation}
 X_m\coloneqq F_{P_1}F_{P_2}\cdots F_{P_m}(X)\quad\text{and}\quad Y_m\coloneqq F_{P_1}F_{P_2}\cdots F_{P_m}(Y)
 \end{equation}
 are contained in the semiring $\AA_+\langle X^{\pm1},Y^{\pm1}\rangle\subset\KK$ of pseudo-positive non-commutative Laurent polynomials.
 \end{maintheorem}
 \begin{introremark}
 When $P_1$ and $P_2$ are monic and of the same degree but $P_1(0)=P_2(0)\ne1$, this result also holds and can be deduced from the Main Theorem by passing to an appropriate algebraic extension of $\kk$, then rescaling all variables. 
 The same is true when the coefficients $p_{1,0},p_{1,d_1},p_{2,0},p_{2,d_2}\ne0$ are arbitrary but satisfy a balancing condition which we leave as an exercise for the reader to work out. 
 In the absence of such a balancing condition the definitions of the polynomials $P_k$ should be adjusted according to the exchange polynomial mutation rules developed by Chekhov and Shapiro~\cite{chekhov-shapiro}.
 Also, since $F_P(X)=QY$ for any polynomial $P$, we have $X_{m+1}=QY_m$ for $m\ge0$; in particular, the claim for the $X_m$ follows from the claim for the $Y_m$.
 \end{introremark}
 \begin{introremark}
 If $d_1d_2\le3$, the Main Theorem can be observed quite explicitly by computing $X_1,X_2,\ldots,X_m$ by hand for 
 \[
 m=\begin{cases} 4 & \text{if $d_1d_2=0$;}\\ 5 & \text{if $d_1d_2=1$;}\\ 6 & \text{if $d_1d_2=2$;}\\ 8 & \text{if $d_1d_2=3$;}\end{cases}
 \]
 and observing in each case that these are given by pseudo-positive non-commutative Laurent polynomials with $X_m=QXQ^{-1}$.
 The combinatorics below can be adapted to these cases, however in everything that follows we assume $d_1d_2\ge4$ as such cases may be treated more uniformly.
 \end{introremark}

 For the following example, observe that the $Y_m$ for $m\ge2$ may alternatively be computed via the following non-commutative analogue of generalized cluster exchange relations:
 \begin{equation}
 \label{eq:non-commutative exchange relation}
 Y_mQY_{m-2}=1+p_{m,1}Y_{m-1}+\cdots+p_{m,d_m-1}Y_{m-1}^{d_m-1}+Y_{m-1}^{d_m}.
 \end{equation}
 \begin{introexample}
 Let $P_1=1+p_{1,1}z+p_{1,2}z^2+z^3$ and $P_2=1+p_{2,1}z+z^2$.
 Then the first few non-commutative generalized cluster variables $Y_m$ are given by:
 \[ Y_1=(1+p_{1,1}Y+p_{1,2}Y^2+Y^3)X^{-1}, \quad Y_2=(1+p_{2,1}Y_1+Y_1^2)Y^{-1}Q^{-1},\]
 \[Y_3=(1+p_{1,2}Y_2+p_{1,1}Y_2^2+Y_2^3)Y_1^{-1}Q^{-1}.\]
 While $Y_2$ is manifestly an element of $\AA_+\langle X^{\pm1},Y^{\pm1}\rangle$, a highly nontrivial cancellation must occur in the expansion of $Y_3$ in order for it to be a pseudo-positive non-commutative Laurent polynomial.
 Such cancellations indeed occur and we obtain the expansion
 \begin{align*}
 Y_3&\Mk =\Mk 
 \Big(XY^{-3}+p_{1,2}(p_{2,1}+Y_1)Y^{-1}+p_{1,1}(p_{2,1}+Y_1)Y^{-2}\\
 &\quad +\Mk p_{1,1}(1+p_{2,1}Y_1+Y_1^2)XY^{-1}X^{-1}(p_{2,1}+Y_1)Y^{-1}\\
 &\quad +\Mk (p_{2,1}+Y_1)Y^{-3}+(1+p_{2,1}Y_1+Y_1^2)XY^{-1}X^{-1}(p_{2,1}+Y_1)Y^{-2}\\
 &\quad +\Mk (1+p_{2,1}Y_1+Y_1^2)XY^{-1}X^{-1}(1+p_{2,1}Y_1+Y_1^2)XY^{-1}X^{-1}(p_{2,1}+Y_1)Y^{-1}\Big)Q^{-1}.
 \end{align*}
 \end{introexample}

 The automorphisms $F_{P_k}$ are generalizations of automorphisms of $\KK$ introduced by Kontsevich~\cite{kontsevich} which are recovered in the \emph{binomial case} when $p_{1,i}=0=p_{2,j}$ for $1\le i\le d_1-1$ and $1\le j\le d_2-1$.
 In this binomial case, Kontsevich conjectured the Laurentness and positivity of the \emph{non-commutative cluster variables} $X_m$ and $Y_m$.
 This terminology is justified by specializing to commutative variables through which we recover the initial cluster mutations in the rank two cluster algebra~\cite{fomin-zelevinsky} associated to the exchange matrix $\left[\begin{array}{cc} 0 & d_2\\ -d_1 & 0\end{array}\right]$ after composing with the transposition of initial cluster variables.
 In the binomial case, Laurentness was established by Usnich~\cite{usnich0} when $d_1=d_2=2$, and by Berenstein and Retakh~\cite{berenstein-retakh} for arbitrary polynomial degrees.
 Positivity in the binomial case was proven by Di Francesco and Kedem~\cite{difrancesco-kedem} when $d_1d_2=4$, by Lee and Schiffler~\cite{lee-schiffler} for $d_1=d_2$, and by the author~\cite{rupel1} for arbitrary polynomial degrees.

 The Laurentness of $X_m$ and $Y_m$ was established by Usnich~\cite{usnich2} in the special case where $P_k=P_1$ for all $k\in\ZZ$.
 We will prove the Main Theorem by providing a combinatorial construction of the elements $Y_m$, called \emph{non-commutative generalized cluster variables}. 
 This combinatorics was studied by the author~\cite{rupel2} to construct greedy bases for (commutative) rank two generalized cluster algebras by building upon the notion of compatible pairs in a maximal Dyck path developed by Lee, Li, and Zelevinsky~\cite{lee-li-zelevinsky} for constructing greedy bases of rank two cluster algebras.

 For $\bfa=(a_1,a_2)\in\ZZ_{\ge0}^2$, let $D\coloneqq D_\bfa$ denote the lattice path in the rectangle $[0,a_1]\times[0,a_2]$ which begins at $(0,0)$ takes unit length East and North steps to end at $(a_1,a_2)$ and is maximal among all such \emph{Dyck paths} that never pass above the main diagonal of the rectangle $[0,a_1]\times[0,a_2]$. 
 In other words, no lattice point of $D$ lies strictly above the main diagonal and any lattice point which lies strictly above $D$ also lies strictly above the main diagonal.
 Label the edges of $D$ as $E=\{1,\ldots,a_1+a_2\}$, where this bijection of ordered sets respects the natural order on edges from $(0,0)$ to $(a_1,a_2)$.
 There is a partition $E=H\sqcup V$, where $H$ (resp. $V$) denotes the set of horizontal (resp. vertical) edges of $D$. 

 For edges $e,e'\in E$, we write $ee'$ for the subpath of $D$ beginning with $e$ traveling North-East and ending with $e'$. 
 By convention, this path will be empty if $e$ is to the North-East of $e'$, while the path $ee$ contains the single edge $e$. 
 Let $\overline{e}e'$ (resp. $e\overline{e}'$) denote the path obtained from $ee'$ by removing the edge $e$ (resp. $e'$). 
 Write $(ee')_H$ (resp. $(ee')_V$) for the set of horizontal (resp. vertical) edges in the path $ee'$.
 We abbreviate $|ee'|_H\coloneqq |(ee')_H|$ and $|ee'|_V\coloneqq |(ee')_V|$. 
 
 \begin{remark}
 In~\cite{lee-li-zelevinsky} and~\cite{rupel2}, the definition for subpaths $ee'$ of $D$ includes a ``wrap-around'' condition whereby $ee'$ is non-empty for $e'<e$, however following~\cite[Remark~2.21]{rupel2} such a condition will not be necessary in our situation and all relevant results quoted from~\cite{rupel2} will be modified accordingly.
 \end{remark}

 \begin{definition}
 \label{def:compatibility}
 A \emph{grading} $\omega:E\to\ZZ_{\ge0}$ (on the edges) of $D$ is called \emph{compatible} if: for every $h\in H$ and $v\in V$, there exists an edge $e$ along the path $hv$ so that at least one of the following holds:
 \begin{align}
 \tag{HGC}\label{eq:hgc} e\ne v\quad&\text{ and }\quad |he|_V=\sum\limits_{h'\in(he)_H} \omega(h');\\
 \tag{VGC}\label{eq:vgc} e\ne h\quad&\text{ and }\quad |ev|_H=\sum\limits_{v'\in(ev)_V} \omega(v').
 \end{align}
 \end{definition}

 Recall that $d_1,d_2\in\ZZ_{\ge0}$ denote the degrees of the exchange polynomials $P_1$ and $P_2$ respectively. 
 We say that a grading $\omega$ of $D$ is \emph{$(d_1,d_2)$-bounded} if $\omega(h)\le d_1$ for all $h\in H$ and $\omega(v)\le d_2$ for all $v\in V$. 
 For the remainder of the paper we will restrict to such bounded gradings $\omega$, though we continue to write $\omega:E\to\ZZ_{\ge0}$ throughout.
 This notion of compatible gradings was introduced in~\cite{rupel2} building upon the notion of compatible subsets of $E$ developed in~\cite{lee-li-zelevinsky} which can be recovered when $\omega(h)\in\{0,d_1\}$ for $h\in H$ and $\omega(v)\in\{0,d_2\}$ for $v\in V$.

 For a $(d_1,d_2)$-bounded grading $\omega$, we associate the non-commutative monomial $\wt_\omega(e)$ to each edge $e\in E$ as follows:
 \begin{equation}
 \label{eq:edge weights}
 \wt_\omega(e)=\begin{cases}
 p_{1,\omega(e)}Y^{\omega(e)}X^{-1} & \text{ if $e\in H$;}\\
 p_{2,d_2-\omega(e)}X^{\omega(e)+1}Y^{-1}X^{-1} & \text{ if $e\in V$.}\\
 \end{cases}
 \end{equation}
 Thus we may associate a non-commutative Laurent monomial to each $(d_1,d_2)$-bounded grading $\omega$ by taking the product of the associated non-commutative weights in the natural order along the path $D$:
 \begin{equation}
 \label{eq:total path weights}
 Y_D(\omega)\coloneqq \wt_\omega(1)\wt_\omega(2)\cdots\wt_\omega(a_1+a_2).
 \end{equation} 
 Define $Y_D\coloneqq \sum\limits_\omega Y_D(\omega)$, where the sum ranges over all $(d_1,d_2)$-bounded compatible gradings $\omega$ of $D$.

 We will mainly be interested in the maximal Dyck paths $D_m\coloneqq D_{\bfa_m}$ for integer vectors $\bfa_m\in\ZZ^2$, $m\ge1$, defined recursively by
 \begin{equation}
 \label{eq:roots recursive}
 \bfa_0=(0,-1),\quad
 \bfa_1=(1,0),\quad
 \bfa_{m-1}+\bfa_{m+1}=\begin{cases}d_2\bfa_m & \text{if $m$ is odd;}\\ d_1\bfa_m & \text{if $m$ is even.}\end{cases}
 \end{equation}
 These vectors are precisely the \emph{almost positive roots} in the root system associated to the Cartan matrix $\left[\begin{array}{cc} 2 & -d_2\\ -d_1 & 2\end{array}\right]$ which describe the denominator vectors of cluster variables.
 The Main Theorem is an immediate consequence of the following combinatorial construction of the non-commutative generalized cluster variables $Y_m$.
